% BHU PYQ -- Manually transcribed & error-corrected
% Paper: BCF-213 : Corporate Taxation
% Exam: B.Com. (Hons.) F.M.M. Semester III Examination, 2023-24

\documentclass[11pt,a4paper]{article}
\usepackage[a4paper,top=2.5cm,bottom=2.5cm,left=2.8cm,right=2.8cm,headheight=14pt]{geometry}
\usepackage{amsmath,amssymb}
\usepackage[T1]{fontenc}
\usepackage[utf8]{inputenc}
\usepackage{lmodern,microtype,fancyhdr,enumitem,hyperref,lastpage,parskip,booktabs}

\hypersetup{
    colorlinks=true,
    urlcolor=black,
    linkcolor=black,
    pdftitle={BCF-213: Corporate Taxation -- BHU B.Com. (Hons.) FMM Sem III 2023-24}
}

\pagestyle{fancy}
\fancyhf{}
\renewcommand{\headrulewidth}{0.4pt}
\renewcommand{\footrulewidth}{0.4pt}
\fancyhead[L]{\small\textit{Banaras Hindu University}}
\fancyhead[R]{\small\textit{BCF-213: Corporate Taxation}}
\fancyfoot[C]{\small Page \thepage\ of \pageref{LastPage}}

\newlist{parts}{enumerate}{2}
\setlist[parts,1]{label=(\alph*),leftmargin=*,itemsep=5pt,topsep=3pt}
\setlist[parts,2]{label=(\roman*),leftmargin=*,itemsep=3pt,topsep=2pt}
\newcommand{\pts}[1]{\hfill\textit{[#1]}}

\begin{document}

\begin{center}
  {\large\bfseries BANARAS HINDU UNIVERSITY}\\[2pt]
  {\normalsize B.Com. (Hons.) F.M.M. --- Semester III Examination, 2023-24}\\[6pt]
  \rule{0.8\linewidth}{0.5pt}\\[6pt]
  {\large\bfseries Paper No. BCF-213: Corporate Taxation}\\[4pt]
  {\normalsize Time: 3 Hours\hfill Maximum Marks: 70}
\end{center}

\rule{\linewidth}{0.4pt}

\smallskip
\noindent\textit{Instructions: Answer all questions. All questions carry equal marks (14 marks each).}

\bigskip

\noindent\textbf{Question 1.}
Define the term Income. Distinguish Gross Total Income and Total Income.\pts{14}

\centerline{\textbf{OR}}

The following are the incomes of Mr.\ 'X' for the previous year 2022-23:\pts{14}
\begin{parts}
  \item Interest on USA Development Bond (1/2 received in India) \hfill Rs.\ 50,000
  \item Income from agriculture in America received there but later on remitted to India \hfill Rs.\ 80,000
  \item Income from property in Canada received outside India \hfill Rs.\ 50,000
  \item Income earned from business in Uganda, which is controlled from Delhi (Rs.\ 25,000 received in India) \hfill Rs.\ 40,000
  \item Dividend paid by domestic company and received outside India \hfill Rs.\ 40,000
  \item Profit from a business in Chennai, which is controlled from London \hfill Rs.\ 50,000
\end{parts}
From the above particulars, ascertain the gross total income of Mr.\ 'X' for the assessment year 2023-24, if he is:
\begin{parts}
  \item a resident (ordinary),
  \item a not ordinarily resident,
  \item a non-resident.
\end{parts}

\medskip\rule{\linewidth}{0.2pt}

\noindent\textbf{Question 2.}
Advise an assessee about the admissibility or otherwise of the claims with regard to the following items, giving reasons:\pts{14}
\begin{parts}
  \item Compensation paid to an employee for premature termination of his services.
  \item Amount spent on successful suit filed against another for infringing the assessee's trade mark.
  \item The penalty paid to customs authority for importing prohibited goods which yielded a large margin of profits.
  \item Travelling expenses of a director who went to Europe for negotiating the purchase of a new heavy machinery which was eventually installed next year.
  \item Cost of erecting a medical annexe to the factory for the emergency treatment of the employees.
  \item Lump-sum consideration paid on 01-07-2022 for acquiring know-how Rs.\ 5,00,000.
  \item Loss due to embezzlement by an employee.
\end{parts}

\centerline{\textbf{OR}}

What deductions are allowed to a businessman in computing profits? Specify the expenses disallowed.\pts{14}

\medskip\rule{\linewidth}{0.2pt}

\noindent\textbf{Question 3.}
Explain and illustrate the term `Book Profit' in relation to the assessment of firms.\pts{14}

\centerline{\textbf{OR}}

`A', `B' and `C' are partners of a firm sharing profit and loss equally. The Profit and Loss Account for the year ended 31st March, 2023 shows a net profit of Rs.\ 99,750 after debiting the following as per partnership deed:\pts{14}
\begin{parts}
  \item Salaries of Rs.\ 20,000 and Rs.\ 15,000 to `A' and `B' respectively.
  \item Bonus to `C' Rs.\ 15,000.
  \item Rs.\ 5,000 for interest on capital to `A' calculated @ 20\%.
  \item Rs.\ 10,000 for rent of the business premises paid to `C'.
  \item Commission of Rs.\ 5,000 to `C'.
\end{parts}
Compute Book Profit and Total Income of the firm for the Assessment Year 2023-24 assuming that it is a professional firm and all are working partners.

\medskip\rule{\linewidth}{0.2pt}

\noindent\textbf{Question 4.}
Shri Bose has estimated the following incomes for the Financial Year 2023-24:\pts{14}
\begin{parts}
  \item Income from house property \hfill Rs.\ 57,00,000
  \item Income from business \hfill Rs.\ 75,00,000
  \item Dividend from Indian Company \hfill Rs.\ 25,000
  \item Interest from Saving Deposit \hfill Rs.\ 1,25,000
\end{parts}
Determine the amount of instalments payable as advance tax (as per the old tax regime).

\centerline{\textbf{OR}}

When does the liability to pay advance tax arise? When is such tax paid and how is it calculated?\pts{14}

\medskip\rule{\linewidth}{0.2pt}

\noindent\textbf{Question 5.}
Describe briefly the procedure of an appeal to the Commissioner (Appeals).\pts{14}

\centerline{\textbf{OR}}

What is the new tax regime theory? How to choose between the old tax regime and the new tax regime?\pts{14}

\vfill
\rule{\linewidth}{0.4pt}
\begin{center}\small
\href{https://bhu-exam.vercel.app/}{Click here to visit our website}\\[4pt]
\href{https://me-aryan.vercel.app/}{Click here to visit our contributor}\\[4pt]
\href{https://quashan-me.vercel.app/}{Click here to visit our contributor}
\end{center}

\end{document}
